\hypertarget{the-planar-surface-emitting-tile-an-alternative-to-the-chip-stack-and-its-collimator}{%
\chapter{The planar surface-emitting tile: an alternative to the chip
stack and its
collimator}\label{planar_tile-the-planar-surface-emitting-tile-an-alternative-to-the-chip-stack-and-its-collimator}}

Module 30. Code: \texttt{scripts/m30\_planar\_tile.py},
\texttt{scripts/m3\_grating\_fdtd.py} (extended with a reflector stack
under the BOX and a metal substrate); results
\texttt{results/m30\_coupling.json},
\texttt{results/m30\_directionality.json},
\texttt{results/m30\_imaging.json}, \texttt{results/m30\_power.json}.
Values are SIM or DERIVED, valid in the regime stated, and should be
reviewed before they are quoted.

\hypertarget{1-result}{%
\section{Result}\label{planar_tile-1-result}}

In the tile proposed in \texttt{docs/scaling\_beyond.md} (Sec. 5), each
engine lies flat. Its emitters continue as weak gratings that radiate
upwards, and the grating aperture is imaged onto the
atoms\textquotesingle{} y axis with demagnification D = tile pitch / 4
um = 1500 (6 mm tiles). No fast-axis collimator is needed, and tile
placement errors shrink by D. We simulated the four quantities the
proposal left open:

\begin{itemize}
\tightlist
\item
  (i) the grating directionality;
\item
  (ii) the y imaging and its angle tolerance;
\item
  (iii) the coupling between neighbouring long emitters;
\item
  (iv) the power per tweezer.
\end{itemize}

(i) \textbf{Directionality is solved by a mirror under the BOX} (2D
FDTD, 30 um weak grating, about 10 deg emission).

\begin{itemize}
\tightlist
\item
  \textbf{Si substrate alone:} at best 0.70 at a 2.65 um BOX, and
  between 0.28 and 0.70 as the BOX varies, so the BOX thickness must be
  controlled to about +/-25 nm.
\item
  \textbf{Bragg mirror, two Si/SiO2 pairs:} 0.965 to 0.985 for a 1.8 to
  2.2 um BOX, except 0.60 at 2.1 um.
\item
  \textbf{Bragg mirror, three pairs:} 0.993 to 0.997, except 0.945 at
  2.1 um.
\item
  \textbf{Perfect metal mirror:} 1.000, but its absorption is not
  modelled.
\item
  \textbf{Deeper etch.} A 0.1 um etch gives the same directionality
  (0.72 / 0.986 / 0.998 / 1.000).
\end{itemize}

(ii) \textbf{The grating must be about 4 mm long, not 2 mm.} The atoms
see the grating\textquotesingle s near field demagnified 1500 times,
through a pupil of object-side NA 0.5 / 1500 = 3.3e-4. A near field that
fills this pupil like the edge emitter\textquotesingle s facet image has
a 1/e\^{}2 radius of about D x 0.97 um = 1.45 mm.

\begin{itemize}
\tightlist
\item
  \textbf{2 mm grating:} truncation leaves crosstalk at 4 um of 3e-3 to
  7e-3 for every profile, 30 to 90 times the edge
  emitter\textquotesingle s 8e-5.
\item
  \textbf{4 mm grating, Gaussian apodization w = 1.7 mm:} 1.2e-4
  crosstalk at a y FWHM of 1.30 um, the shortest design that meets it.
\item
  \textbf{Angle tolerance of that design:} 5.2e-5 in sine for peak
  \textgreater= 0.97, 3.0e-5 for 0.99, and 3.6e-5 for crosstalk
  \textless= 2e-4. This confirms the 30 urad estimate. Object defocus is
  irrelevant at this NA: 1.6 mm changes the peak by less than 1e-7 (2e-4
  um at the atoms).
\end{itemize}

(iii) \textbf{Neighbouring emitters couple, and this sets the y angle
per tweezer.}

\begin{itemize}
\tightlist
\item
  \textbf{Mechanism.} On a 4 mm grating the emitters form a coupled
  array. A tweezer at x\_a is close to one Bloch mode, whose effective
  index, and so whose y emission sine, depends on x\_a. The swing across
  the field is 4 kappa / k0.
\item
  \textbf{Uniform array.}

  \begin{itemize}
  \tightlist
  \item
    At the D1 pitch of 1.5 um the swing is 4.7e-3, 90 times the
    tolerance.
  \item
    Pitch 2.0 um: 1.6e-4. Pitch 2.3 um: 2.2e-5, which meets the
    tolerance; 2.2 um is about the limit.
  \item
    A wider pitch narrows the grating-lobe-free field to lambda / (2p):
    0.65 of D1 at 2.3 um, so each engine serves 35\% fewer sites, or
    needs 1.5 times more ports.
  \end{itemize}
\item
  \textbf{Alternating widths.} Widths of 1.0 / 1.4, 0.9 / 1.5 or 0.8 /
  1.6 um at 1.5 um pitch cut the swing to 2.1e-4, 1.4e-4 and 1.0e-4: 22
  to 45 times, but still 2 to 4 times over the tolerance. The two
  emitter types also radiate at different angles unless their grating
  periods differ.
\item
  \textbf{Method check.} These come from 8-guide array supermodes. The
  two-guide extraction fails for mismatched guides (Sec. 3); for uniform
  arrays it agrees with the 8-guide result (4.65e-3 against 4.66e-3).
\end{itemize}

(iv) \textbf{Power per 1.2 mK tweezer:} 46.5 mW with a three-pair Bragg
mirror (47.1 mW with two pairs, 66 mW with the substrate alone), against
57.7 mW for the edge emitter with the reflective collimator. This counts
the y axis only; the x axis is that of the edge emitter, at the pitch D1
uses.

\textbf{At the wider pitch the coupling needs, the gain is lost
(provisional).} The 35\% loss of x field at 2.2 um pitch may be
recoverable by changing the x-relay magnification, since the number of
resolvable spots equals the number of emitters at any pitch. The price
is the fill factor (1.2 um emitters: 0.80 at 1.5 um, 0.55 at 2.2 um).
The fraction of one emitter\textquotesingle s far-field power inside the
grating-lobe-free range \textbar s\textbar{} \textless= lambda / (2p)
falls from 0.71 at 1.5 um to 0.54 at 2.2 um (0.51 at 2.3 um), a relative
efficiency of 0.76 (DERIVED from the femwell emitter mode; the x relay
is not simulated). The tile at 2.2 um would then need about 61 mW per
tweezer, no better than the edge emitter.

Hence the tile removes the stack, the collimator and its +/-0.1 um glue
tolerance. It is about 20\% more efficient only at the D1 pitch, where
coupling rules it out; at the pitch coupling allows it is about equal to
the edge emitter. It brings its own costs:

\begin{itemize}
\tightlist
\item
  a 4 mm apodized grating per emitter, with the emission rate controlled
  along it;
\item
  a Bragg mirror under the BOX;
\item
  an emitter pitch of about 2.2 um, or a phase-mismatched emitter design
  that has not been found;
\item
  a 96 mm y-relay object field for 16 tiles. The coupling problem is the
  open item. It does not arise for the edge emitter, whose emitters end
  at the facet.
\end{itemize}

\hypertarget{2-directionality-meep}{%
\section{Directionality (MEEP)}\label{planar_tile-2-directionality-meep}}

2D FDTD of the partially etched SiN grating:

\begin{itemize}
\tightlist
\item
  400 nm film, 50 nm etch, duty 0.5, 30 um long;
\item
  top oxide 2 um;
\item
  40 px/um;
\item
  normalised to a straight-guide reference. The original cases are
  reproduced exactly (0.715 at 0.1 um etch, 3 um BOX).
\end{itemize}

{\small
\begin{longtable}[]{@{}>{\raggedright\arraybackslash}p{0.295\dimexpr(\linewidth-8\tabcolsep)\relax}>{\raggedright\arraybackslash}p{0.098\dimexpr(\linewidth-8\tabcolsep)\relax}>{\raggedright\arraybackslash}p{0.295\dimexpr(\linewidth-8\tabcolsep)\relax}>{\raggedright\arraybackslash}p{0.311\dimexpr(\linewidth-8\tabcolsep)\relax}@{}}
\toprule\noalign{}
reflector & best BOX & directionality (0.05 / 0.1 um etch) & range over
the BOX values simulated \\
\midrule\noalign{}
\endhead
\bottomrule\noalign{}
\endlastfoot
Si substrate & 2.65 um & 0.702 / 0.722 & 0.28 to 0.70 (2.0 to 2.8 um, 50
nm steps) \\
Bragg, 2 pairs (Si 75 nm / SiO2 184 nm) & 1.9 um & 0.985 / 0.986 & 0.60
to 0.985 (1.8 to 2.2 um) \\
Bragg, 3 pairs & 1.9 um & 0.997 / 0.998 & 0.945 to 0.997 \\
perfect metal & 1.8 to 2.2 um & 1.000 / 1.000 & 1.000 \\
\end{longtable}
}

The BOX also sets the emission rate. The 1/e power length at 50 nm etch
runs from 115 um to 960 um with the metal mirror, and from 124 um to 2.7
mm with the Bragg mirror. An apodized 4 mm grating therefore needs both
its etch (about 15 nm for a decay length of order 1.5 mm, scaling as
etch\^{}-2 from these data) and the BOX controlled along its length.
Back-reflection into the guide is below 1e-3.

\hypertarget{3-coupling-of-long-emitters}{%
\section{Coupling of long
emitters}\label{planar_tile-3-coupling-of-long-emitters}}

Supermodes of 1.2 x 0.4 um SiN emitters at 1064.4 nm (femwell). For a
uniform array the two-guide splitting gives kappa, and the Bloch band
spans 4 kappa / k0 = 2 x (splitting):

{\small
\begin{longtable}[]{@{}>{\raggedright\arraybackslash}p{0.124\dimexpr(\linewidth-10\tabcolsep)\relax}>{\raggedright\arraybackslash}p{0.221\dimexpr(\linewidth-10\tabcolsep)\relax}>{\raggedright\arraybackslash}p{0.195\dimexpr(\linewidth-10\tabcolsep)\relax}>{\raggedright\arraybackslash}p{0.159\dimexpr(\linewidth-10\tabcolsep)\relax}>{\raggedright\arraybackslash}p{0.301\dimexpr(\linewidth-10\tabcolsep)\relax}@{}}
\toprule\noalign{}
pitch (gap) & splitting & coupling length & band swing in sine & field
lambda / 2p (fraction of D1) \\
\midrule\noalign{}
\endhead
\bottomrule\noalign{}
\endlastfoot
1.5 um (0.3) & 2.3e-3 & 0.23 mm & 4.7e-3 & 0.355 (1.00) \\
1.7 um (0.5) & 6.0e-4 & 0.89 mm & 1.2e-3 & 0.313 (0.88) \\
2.0 um (0.8) & 8.0e-5 & 6.7 mm & 1.6e-4 & 0.266 (0.75) \\
2.3 um (1.1) & 1.1e-5 & 48 mm & 2.2e-5 & 0.231 (0.65) \\
2.6 um (1.4) & 1.5e-6 & 0.35 m & 3.1e-6 & 0.205 (0.58) \\
\end{longtable}
}

For alternating widths the two-guide extraction fails. The splitting at
1.5 um pitch falls below the index difference of the isolated guides (by
5e-5 to 8e-4), because each guide\textquotesingle s index is pulled by
its neighbour\textquotesingle s dielectric, and unequally. We therefore
read the band width from the spread of the four upper-band supermodes of
an 8-guide array. For a chain of m guides, n\_j = n0 + c cos(j pi / (m +
1)), and the infinite band spans 2c. The same method on uniform 8-guide
arrays reproduces the two-guide values (4.65e-3 at 1.5 um, 1.60e-4 at
2.0 um).

\hypertarget{4-imaging-and-power}{%
\section{Imaging and power}\label{planar_tile-4-imaging-and-power}}

1D scalar Fourier optics:

\begin{itemize}
\tightlist
\item
  near field A(u) over the grating, demagnified by D = 1500;
\item
  object-side pupil \textbar s\textbar{} \textless= 3.3e-4;
\item
  normalised so that the line integrates to the fraction of the emitted
  power that passes the pupil, as in the wave model of module 13.
\end{itemize}

{\small
\begin{longtable}[]{@{}>{\raggedright\arraybackslash}p{0.383\dimexpr(\linewidth-8\tabcolsep)\relax}>{\raggedright\arraybackslash}p{0.266\dimexpr(\linewidth-8\tabcolsep)\relax}>{\raggedright\arraybackslash}p{0.234\dimexpr(\linewidth-8\tabcolsep)\relax}>{\raggedright\arraybackslash}p{0.117\dimexpr(\linewidth-8\tabcolsep)\relax}@{}}
\toprule\noalign{}
profile & crosstalk at 4 um & peak per emitted power & y FWHM \\
\midrule\noalign{}
\endhead
\bottomrule\noalign{}
\endlastfoot
2 mm, uniform weak grating (Le 1 mm) & 3.3e-3 & 0.73 & 1.03 um \\
2 mm, Gaussian w 1.2 mm & 3.6e-3 & 0.80 & 1.02 um \\
4 mm, Gaussian w 1.45 mm & 3.9e-4 & 0.78 & 1.22 um \\
4 mm, Gaussian w 1.7 mm (adopted) & 1.2e-4 & 0.73 & 1.30 um \\
5 mm, Gaussian w 1.7 mm & 4e-5 & 0.69 & 1.38 um \\
\end{longtable}
}

Tolerances of the adopted profile:

\begin{itemize}
\tightlist
\item
  \textbf{Emission-angle error:} 20 / 40 / 60 / 100 urad give peaks of
  0.996 / 0.983 / 0.962 / 0.896 and crosstalk of 1.4e-4 / 2.2e-4 /
  3.9e-4 / 1.1e-3.
\item
  \textbf{Decentre:} the tile\textquotesingle s y placement moves the
  row by 1/1500 of the error.
\end{itemize}

\textbf{Power.} The edge chain delivers a line peak of 0.59 per unit
guided power, from the facet transmission into glass, the pupil fill and
a silvered M1 with a TIR M2 (m14, m27). The tile delivers 0.733 x
directionality. The power per tweezer scales inversely, which gives the
figures of Sec. 1(iv).

\textbf{Layout.} 16 tiles at 6 mm make a 96 mm y-relay object field at
NA 3.3e-4. Its Lagrange invariant matches the atoms (48 mm x 3.3e-4 = 16
rows x 4 um / 2 x 0.5). A 4 mm grating takes two thirds of the 6 mm tile
pitch along the grating, so the engine (FPR slab of about 2 mm) has to
sit beside the emitter array in x rather than in line.

\hypertarget{5-limits}{%
\section{Limits}\label{planar_tile-5-limits}}

\begin{itemize}
\tightlist
\item
  \textbf{Dimensionality.} The FDTD is 2D (infinitely wide emitters)
  with a uniform 30 um grating. Directionality is taken as independent
  of length and apodization. Metal absorption is not modelled.
\item
  \textbf{Coupling.} The coupling is computed for straight guides, not
  for the etched grating region, which changes the modes slightly. The
  tolerance applies to the swing across the whole field; partial
  compensation by trimming each emitter\textquotesingle s grating period
  was not studied.
\item
  \textbf{Imaging.} The imaging model is scalar and 1D in y; the x axis
  and the circular pupil are not combined here.
\item
  \textbf{Engine.} The engine itself (FPR, actuators) is unchanged and
  not re-simulated at a wider emitter pitch.
\end{itemize}

\hypertarget{verification-notes}{%
\section{Verification notes}\label{planar_tile-verification-notes}}

{\footnotesize
\begin{longtable}[]{@{}>{\raggedright\arraybackslash}p{0.309\dimexpr(\linewidth-8\tabcolsep)\relax}>{\raggedright\arraybackslash}p{0.117\dimexpr(\linewidth-8\tabcolsep)\relax}>{\raggedright\arraybackslash}p{0.101\dimexpr(\linewidth-8\tabcolsep)\relax}>{\raggedright\arraybackslash}p{0.473\dimexpr(\linewidth-8\tabcolsep)\relax}@{}}
\toprule\noalign{}
value & location & type & pointer \\
\midrule\noalign{}
\endhead
\bottomrule\noalign{}
\endlastfoot
directionality, emission rate, back-reflection & Sec. 1, 2 & SIM &
results/m30\_directionality.json (2D FDTD; review) \\
coupling, band swing, field fraction & Sec. 1, 3 & SIM &
results/m30\_coupling.json (review) \\
imaging profiles, angle and focus tolerance & Sec. 1, 4 & SIM &
results/m30\_imaging.json \\
power per tweezer & Sec. 1, 4 & DERIVED & results/m30\_power.json \\
capture fraction vs pitch, relay-magnification alternative & Sec. 1 &
DERIVED & results/m30\_power.json (x\_relay\_alternative;
provisional) \\
tile pitch 6 mm, D = 1500, grating length choices & Sec. 1 & CHOICE &
docs/scaling\_beyond.md Sec. 5 \\
n(Si) 3.55, n(SiO2) 1.4496 at 1064.4 nm & Sec. 2 & SOURCE / CHOICE &
picsim/params.py, picsim/materials.py \\
edge chain 57.7 mW & Sec. 1 & SIM & results/m27\_wave.json \\
\end{longtable}
}

Open placeholders: {[}REF?{]} fabrication of SiN gratings with a Si/SiO2
Bragg mirror under the BOX; {[}VALUE?{]} absorption of a practical metal
mirror at 1064 nm under oxide.
